% ===========================================================================
% OWNER: JOINT - lead M3 (Zarpash Nasim); each member adds their module's remaining work by PR
% GRADED BY: Form 3 #3 (feasibility), Form 2 #2
% DRAW FROM: docs/00-Meta/roadmap.md, docs/08-MidEval/iteration-plan.md, docs/05-Testing/evaluation-plan.md, docs/07-Defense/decisions/cdr-gated-module.md
% Rules: mid-eval-report/CLAUDE.md. Do not edit a file you do not own.
% ===========================================================================
\chapter{Remaining FYP Plan}

\section{Remaining Scope}
Scope was baselined at the proposal and has not been reopened; what remains is what was already
committed but not yet built, plus two explicitly gated tiers (\path{docs/02-Requirements/
scope-boundary.md}). \textbf{Committed, not yet built:} the filesystem resource-event family
(static fields and system properties are built; filesystem is not); the cost-aware adaptive
search policy (ranked one-by-one ships now); rootless container sandboxing and the rest of the
security posture (Chapter 4's Table~\ref{ch4:tab:fr} and Appendix B's Table~\ref{appb:tab:security}
list exactly what is and is not built today); the test--resource graph beyond pairwise
edges; and NFR-2's clean-machine deployment check by a third party. \textbf{Gated, conditional
on the committed core:} Certificate-Driven Remediation (CDR) --- enabled only after the
committed core reaches its Final-1 exit criteria and ranking/explanation quality clears an
acceptance threshold on the development benchmark (\path{docs/07-Defense/decisions/
cdr-gated-module.md}) --- and the other optional tier (JUnit 5, Gradle, GitHub Checks,
issue-tracker export, a history dashboard, database/network resource adapters), each gated the
same way. If the team is behind at Final-1, the recorded cut order is CDR first, then
filesystem evidence (reduced to an opaque event), then the adaptive policy (ranked one-by-one
stays default); \textbf{the committed core end-to-end path is cut last, because it is the
project.}

\section{Remaining Modules and Features}
\textbf{Evaluation (M3, confident --- this is my own module).} Integrate at least one published
baseline (iFixFlakies minimiser, or RankFO) so claim B11 can move from ``partially'' to
measured; run the ablations already designed but not executed (resource-only, order-only,
static-only, no-cost-term, no-verification-reserve, the depth-sweep reporting both rank and
overlap-set size); build the held-out evaluation split (development / validation / untouched
projects, never split randomly by test); run the five user-acceptance tasks with the named
participant (SH1) and a substitute plan for the other two.

\todo{M1: please confirm --- from \path{docs/00-Meta/roadmap.md}'s FYP-I-end and FYP-II-early
phases, your remaining work is: the filesystem resource family; the test--resource graph beyond
pairwise edges; dynamic/runtime evidence collection if static evidence proves insufficient
(Iteration 2, ADR-001); the depth matrix and over-approximation measurement at full scale.}

\todo{M2: please confirm --- from the same roadmap phases, your remaining work is: the
cost-aware adaptive search policy (ranked one-by-one is the Iteration-1 default); container
sandboxing and the full security posture (rootless, network-disabled, resource limits,
redaction, audit); the verification-and-abstention hardening the roadmap calls out for
FYP-II-early; CI/CD packaging and the clean-install deployment check (NFR-2).}

\section{Remaining Iterations}
Table~\ref{ch10:tab:iterations} maps the roadmap's own iteration plan
(\path{docs/00-Meta/roadmap.md}) to what each iteration must produce as exit evidence, not just
what work happens --- the roadmap itself insists a timeline answer is about evidence at each
gate, not a list of activities.

\begin{table}[H]\centering
\caption{Remaining iterations and their exit evidence}\label{ch10:tab:iterations}
\begin{tabular}{p{1.6cm}p{2.2cm}p{9.3cm}}
\toprule
\textbf{Iteration} & \textbf{Months} & \textbf{Exit evidence (not activities)} \\
\midrule
2 & Nov -- Jan & Resource normalisation and the full graph; dynamic collection if static proves
  insufficient; the adaptive policy and ablation harness; deletion minimisation (already landed
  early, W10/ADR-007) \\
3 & Feb -- Apr & Failure-signature and repeated-run verification hardening; the controlled
  fixtures extended and a reproducibility bundle; the certificate's full evidence path;
  CLI/CI and container security; the non-perturbation control as an automated check \\
4 & May -- Jun & End-to-end integration on real projects; user acceptance tasks; final testing,
  deployment and handover \\
\bottomrule
\end{tabular}
\end{table}

\section{Implementation Plan}
My own near-term order: (1) wire one published baseline into \path{eval/baseline.py}'s existing
\texttt{OrderRunner} interface, reusing the class-first shuffle discipline already built and
tested; (2) extend the random-order comparison from the seeded fixture to the \texttt{idoft}
benchmark cases already listed (metadata-only today) in \path{eval/benchmark/manifest.json};
(3) build the held-out split once a second real project is cloned, per the leakage rule (a fix
commit used as an algorithm input cannot also be the oracle that judges it); (4) run the
acceptance tasks last, since they need a stable, demonstrable pipeline rather than one still
changing week to week.

\todo{M1/M2: please add your own near-term implementation order here, or confirm mine does not
block yours --- the evaluation work above can start independently of the filesystem family or
the adaptive policy, but the held-out real-project work benefits from both being closer to done.}

\section{Testing and Validation Plan}
The evaluation plan (\path{docs/05-Testing/evaluation-plan.md}) already defines what remains to
be measured, not just tested: eight research questions (RQ1--RQ6 committed, RQ7--RQ8 gated with
CDR), six baselines (random, one-by-one, iFixFlakies, RankFO/RankFL, Takuan, an oracle ceiling),
and a pre-registered failure condition --- a Verified certificate later shown wrong is the single
worst result the project can produce, more serious than any baseline underperforming. Testing
itself (does the implementation meet its requirements) stays continuous, as in Chapter 7;
evaluation (is the approach better under budget) and acceptance (would the intended users use
it) are what remain, and neither is started yet, stated plainly rather than implied as done.

\section{Deployment Plan}
\todo{M2: please confirm --- NFR-2 commits to: a third party installs on clean Linux from
documentation, runs one frozen benchmark case, obtains the same certificate class, replays the
returned order, and removes all created containers and workspaces, with the two-phase network
model (a networked dependency-fetch stage, then a disabled-network test-execution stage as a
separate container invocation) keeping test code itself never network-reachable even though the
build was. This is explicitly a Final-2 exit requirement, not attempted in Iteration 1.}

\section{Tentative Timeline}
Table~\ref{ch10:tab:gates} gives the remaining gates from \path{docs/00-Meta/roadmap.md}, each
with what must be true to pass it --- the roadmap's own framing, not a restatement of the
iteration table above.

\begin{table}[H]\centering
\caption{Remaining gates}\label{ch10:tab:gates}
\begin{tabular}{p{2.6cm}p{10.5cm}}
\toprule
\textbf{Gate} & \textbf{What must be true} \\
\midrule
Final-1 & Working primary scenario end-to-end; integration proven; baseline measured; first
  deployment attempted \\
Mid-2 & Committed scope integrated; system and non-functional testing complete; evaluation
  underway; failure/edge cases exercised; release/deployment rehearsed \\
Job fair (mid-April) & Stable enough for professional repeated demonstration --- the roadmap is
  explicit that essential development should not be left to after this point \\
\bottomrule
\end{tabular}
\end{table}

\section{Expected Final Deliverables}
Per the roadmap's FYP-II-final exit evidence: frozen held-out results, a recorded acceptance
decision from the user tasks, a deployed release passing the NFR-2 clean-install check, and a
one-command reproducibility bundle. CDR ships only as V1 first, and only if its own gates (ranking
quality, explanation correctness, the committed core's Final-1 exit criteria) are met --- its
absence does not block any of the other deliverables, by design (``remove CDR and the product
still delivers its contribution,''
\path{docs/07-Defense/decisions/cdr-gated-module.md}).
